
\subsubsection{6.1 Key Findings}

The two MUST_WORK sub-hypotheses (H-E1 and H-M1) were satisfied under synthetic pipeline validation. The strong partial correlations across all five trustworthiness metrics, and the dominance of a single factor explaining 72.1\% of shared variance, are consistent with the hypothesized latent structure. The survival fraction of 0.943 indicates that the ECE–TruthfulQA\% association is not attributable to shared MMLU capability: controlling for MMLU reduces the correlation by 5.7\%.

These results are consistent with the hypothesis that epistemic reliability constitutes a dimension largely orthogonal to raw capability. If this pattern were to replicate with real LLM evaluations, it would imply that MMLU-based model screening is not informative about this safety dimension. However, this implication rests entirely on real-data replication and cannot be derived from the synthetic pipeline validation reported here.

One mechanistic interpretation is that training procedure — specifically, the use of RLHF versus supervised fine-tuning — influences epistemic reliability largely independently of capability, which is driven by pretraining scale and data. Within-regime stratified analysis (FW4) would directly test this interpretation.

The single-factor structure, explaining 72.1\% of shared variance, suggests the trustworthiness metric space is low-dimensional in the synthetic data. If confirmed with real evaluations, a compact 2–3 metric battery might efficiently capture epistemic reliability for deployment screening.

\subsubsection{6.2 The ΔAUC Null Result}

The ΔAUC result (0.051, CI [−0.194, 0.449]) does not support the pre-specified threshold of ΔAUC ≥ 0.10. The CI width of 0.643 is a consequence of the combination of N=30, a binary LOO outcome, and a modest underlying effect. This result is uninformative about whether a true incremental advantage exists: the CI is consistent with effect sizes ranging from slightly negative to strongly positive. A minimum sample size of N≥100 and a continuous outcome measure are necessary for a well-powered test of this prediction (FW3).

\subsubsection{6.3 Limitations}

\textbf{L1 — Synthetic data (critical).} All quantitative results reflect a parametric data generator, not real LLM evaluations. The generator has pre-wired latent factor structure, which means the observed correlations and factor solution are properties of the generator rather than empirical discoveries about LLMs. All claims about LLM behavior require real-data replication (FW1) before they can be made.

\textbf{L2 — N=30 underpowered for ΔAUC.} The CI width ([−0.194, 0.449]) renders the ΔAUC result uninterpretable as evidence for or against incremental predictive value.

\textbf{L3 — H-M3 not executed.} The mechanistic pathway from calibration to adversarial robustness via decision-surface smoothness under Gaussian perturbation was pre-registered but not executed. The proposed mediation test (Jonckheere-Terpstra dose-response + bootstrap mediation at ε ∈ {0.005, 0.01, 0.02}) remains an open empirical question.

\textbf{L4 — Decoding invariance not tested.} Tucker's congruence was assessed within the greedy decoding regime only. Cross-condition stability under T=0.7 stochastic decoding is pre-registered but not executed; T=0.7 data were unavailable.

\textbf{L5 — Observational design.} All results are cross-sectional correlations. Causal language about the relationship between calibration and hallucination or adversarial robustness is not warranted.

\textbf{L6 — Training regime metadata unverified.} Model card labels for RLHF, SFT, or base-only training were not independently confirmed.

\textbf{L7 — Within-family non-independence.} Models from the same family (e.g., multiple LLaMA-2 or Mistral parameter sizes) share pretraining data, architecture, and alignment procedure, violating the independence assumption of classical psychometric analysis. With 8 families across 30 models, intra-family correlation may inflate estimates of factor stability. Family-stratified sensitivity analysis (FW4) would directly assess whether the latent structure persists within individual families.

\subsubsection{6.4 Future Work}

\textbf{FW1 — Real-data replication (immediate priority).} Execute the pre-registered pipeline (main.py + lm-evaluation-harness v0.4.x) on the 30-model population. The implementation is complete.

\textbf{FW2 — H-M3 embedding perturbation probe.} Gaussian noise injection at ε ∈ {0.005, 0.01, 0.02}, Jonckheere-Terpstra dose-response test, bootstrap mediation analysis.

\textbf{FW3 — Scale to N≥100.} Required for a well-powered ΔAUC test.

\textbf{FW4 — Training regime stratified analysis.} Test whether the epistemic reliability factor persists within RLHF versus SFT strata, and within individual model families.

\textbf{FW5 — Extension to post-2024 and larger models.} Test replication in models exceeding 70B parameters and released after the 2024-01 cutoff.

\textbf{FW6 — Compact screening battery.} If FW1 confirms the latent structure, develop a minimal 2–3 metric proxy for practical deployment screening.

\subsubsection{6.5 Broader Impact}

The primary value of this work, pending real-data replication, is methodological: it demonstrates a psychometric analysis framework for LLM trustworthiness that can be applied to any population of models with standardized benchmark scores. The primary risk is misinterpretation of synthetic-data results as empirical findings about real LLMs. This limitation is stated explicitly in the abstract, the methodology section, the results section, and the limitations section.

